% Backcast presentation. Build from this folder with `latexmk` (LuaLaTeX, see latexmkrc).
\documentclass[aspectratio=169]{beamer}
\usetheme{dtumocha}

% Slide figures: collected ones in figures/, analysis plots straight from the repo.
\graphicspath{{figures/}{../../analysis/plots/}}

\title{Backcasting Balmorel}
\subtitle{Validating the pan-European model against ENTSO-E data for 2023}
\author{Mathias Berg Rosendal}
\institute{DTU Management, Technical University of Denmark}
\date{\today}

\begin{document}

\maketitle

% ---------------------------------------------------------------------------
\begin{frame}{Template: text}
  \begin{itemize}
    \item A bullet point with \alert{alerted text} and \textcolor{ctpGreen}{example text}
    \item Second level
      \begin{itemize}
        \item Sub-bullet in a muted colour
      \end{itemize}
    \item Numbers in text mode keep the Inter font: 2,719.5 vs.\ 2,718.4\,TWh
  \end{itemize}

  \begin{block}{Takeaway}
    Several of the biggest gaps were in the ruler, not the model.
  \end{block}
\end{frame}

% ---------------------------------------------------------------------------
\begin{frame}{Template: full-width figure}
  \centering
  \includegraphics[width=\textwidth, height=0.78\textheight, keepaspectratio]%
    {backcast_ablation_tornado_F2023}
  \source{Leave-one-out runs of the full-year 2023 backcast; each bar is ablation $-$ reference.}
\end{frame}

% ---------------------------------------------------------------------------
\begin{frame}{Template: figure and text}
  \begin{columns}[c]
    \begin{column}{0.58\textwidth}
      \centering
      \includegraphics[width=\linewidth, height=0.75\textheight, keepaspectratio]%
        {backcast_ablation_tornado_generation_F2023}
    \end{column}
    \begin{column}{0.38\textwidth}
      \begin{itemize}
        \item Without nuclear availability: nuclear \alert{+161}\,TWh
        \item Without hydro FLH: hydro \alert{+84}\,TWh
        \item Gas absorbs most of it
      \end{itemize}
    \end{column}
  \end{columns}
\end{frame}

\end{document}
